\cotefile{18}
\scanpage{33}
Structure \inserted{de \unsure{vectoriel|naturel}} complexe sur $DR(M)$ ; mais on notera que
la bigraduation déduite par transport de structure
de celle de Hodge ne \uline{dépend pas de $u$}
(donc, si on remplace $u(\lambda)$ par $\overline{u(\lambda)}$, elle n'est \uline{pas}
remplacée par la symétrique !).

\inserted{Néanmoins, si on regarde,
pour un motif donné $M$,
le motif $\unsure{{}^{\sigma}M_{\underline{\mathbb C}}|{}^{\sigma}\!M_{\underline C}}$ ($\underline{\mathbb C}$ la clôt.
normale de $K$) alors $DR({}^{\sigma}M_{\underline{\mathbb C}})\xrightarrow{\sim}DR(M)_{\underline{\mathbb C}}$, les structures rationnelles \inserted{et \unsure{filtrations|filtration}} sont les \unsure{mêmes},
\unsure{Ainsi|Alors|mais} \deleted{la filtration} \unsure{composant|composé} avec $\sigma$, on trouve un \uline{$K$-isom}
$DR(M_{\underline{\mathbb C}})\simeq DR(M)$, conservant la $\unsure{Q|\mathbb Q}$-structure, et \uuline{renversant} filtration.}

\deleted{Supp.} Soit maintenant $K_0=\mathbb R$, et considérons les
motifs \inserted{$M$} \uline{sur} $\mathbb R$. Choisissons une clôture alg.
$K$ de $K_0=\mathbb R$. Alors $DR(M)_K=DR(M_K)$
est muni de \uline{deux façons} d'une $\mathbb R$-structure :
du fait qu'il est déduit de $DR(M)$ par $\mathbb R\unsure{-|\to}K$,
et puis : la construction donnée plus haut.
Soit $\sigma$ l'action de conjugaison complexe définie
plus haut, (N.B. $\sigma$ \uline{ne respecte pas} la \deleted{\unsure{bigr.}}\inserted{filtration})
$\tau$ celle définie par $DR(M)_K\simeq DR(M_K)$ (donc
\margin{$X(K)$}$\sigma$ respecte la filtration) -- $\sigma$ et $\tau$ sont
\uline{antilinéaires}. Alors $\sigma$ et $\tau$ commutent,
en d'autres termes $DR(M_K)^\sigma$ $\left(=H^*(X(K),\mathbb R)\ \text{si }M=\mathrm{Mot}^*(X)\right)$
est stable par $\tau$ : en effet, $\tau$ n'est autre
que l'action de $H^*(X(K),\mathbb R)$ définie par
l'action $\tau_X$ de $X(K)$ ! Alors $\sigma\tau=\tau\sigma=f^{\deleted{\unsure{*}}}_{\mathbb R}$
est l'action de $H^*(X(K),K)\simeq DR(\mathrm{Mot}^*(X))$
% archive sheet number: 32
